\subsection{Validation}

The Validation stage verifies the deployment descriptor $D_{\mathrm{deploy}}$
before service deployment. To provide an independent check, the validator does
not reuse the evaluation procedure of the planner: every constraint is
re-derived from the original contracts, namely the offloadability profile $O$,
the resource profile $R$, the infrastructure model $\mathcal{I}$ and the SLA
policy $P_{\mathrm{SLA}}$, and the descriptor is treated as untrusted input.
A planner error therefore cannot validate itself.

\paragraph{Constraint set.}
The set of constraints $\mathcal{C}$ comprises two groups. Per-application
constraints verify that
(V1) every execution unit is deployed on known nodes;
(V2) non-offloadable units run only on the node bound to their anchor;
(V3) initialization units are deployed on exactly the nodes hosting the units
they support;
(V4) offloaded units run on eligible hosts, i.e., not on microcontrollers nor
on local devices of other users or applications;
(V5) GPU and external-service requirements are met by every hosting node;
(V6) proximity conditions hold between the nodes hosting related units;
(V7) replication follows the policy: one replica per node, the replica count
does not exceed the SLA maximum, stateful units are not replicated, and the
replicas required by the SLA are deployed;
(V8) movable units meet the availability target;
(V9) offloaded units have no unverified assumptions;
(V10) units connected by a hardware edge share a node;
(V11) the link loads reported in the descriptor match an independent
recomputation; and
(V12) every flow has a network path.
Joint constraints are evaluated with all applications deployed together on the
shared infrastructure: (J1) the resources reserved on each node and (J2) the
traffic on each link remain below the utilization target $\theta$, and (J3) the
start-up peak of every initialization unit fits in the capacity left on its
node. Resource reservations and traffic are recomputed with the same
provenance-based safety margins and replication semantics (load-split,
mirrored, partitioned) as those defined in $P_{\mathrm{SLA}}$.

\paragraph{Severity.}
Each constraint is classified as an \emph{error} or a \emph{warning}. Errors
correspond to hard constraints and make the descriptor invalid. Warnings
correspond to soft constraints of the policy, such as availability in soft
mode or unverified assumptions when strict assumptions are disabled, and to
consistency checks; they are reported but do not reject the plan. Switching the
policy to hard availability or strict assumptions promotes the corresponding
warnings to errors. The validation result is therefore
\begin{equation}
V_{\mathrm{val}} =
\begin{cases}
\mathrm{Valid}, & \text{if } \forall c\in\mathcal{C}_{\mathrm{err}}:\ D_{\mathrm{deploy}}\models c,\\
\mathrm{Invalid}, & \text{otherwise},
\end{cases}
\end{equation}
where $\mathcal{C}_{\mathrm{err}}\subseteq\mathcal{C}$ is the set of hard
constraints.

\paragraph{Feedback.}
A valid descriptor is accepted as the initial deployment plan. When violations
are detected, they are returned to the Deployment Planner as a feedback
descriptor listing, for each violating unit, the nodes on which it must not be
placed, together with the violated constraints. The planner removes these nodes
from the candidate sets and plans again, and the new descriptor is validated in
turn until it is valid or no feasible plan remains.

\paragraph{Fault injection.}
To assess the validator itself, correct descriptors are deliberately corrupted
with representative planning errors (e.g., moving an anchored unit, placing a
GPU unit on a node without GPU, separating an initialization unit from its
consumers, replicating a stateful unit, removing an SLA-required replica,
co-locating replicas, or offloading a unit whose traffic saturates a link).
The detection rate, i.e., the fraction of corrupted descriptors declared
invalid, measures the coverage of the constraint set.

The procedure is summarized in Algorithm~\ref{alg:validation}.

\begin{algorithm}[ht]
\caption{Deployment Plan Validation}
\label{alg:validation}
\begin{algorithmic}[1]
\Require Deployment descriptors $\{D_{\mathrm{deploy}}^a\}$ of all applications
\Require Contracts $O$, $R$, infrastructure $\mathcal{I}$, SLA policy $P_{\mathrm{SLA}}$
\Ensure Validation results $\{V_{\mathrm{val}}^a\}$, feedback $F$
\State $\mathcal{C} \leftarrow \mathrm{DeriveConstraints}(O, R, \mathcal{I}, P_{\mathrm{SLA}})$ \Comment{independent of the planner}
\State $F \leftarrow \emptyset$
\For{each application $a$}
  \State $V_{\mathrm{val}}^a \leftarrow \mathrm{Valid}$
  \For{each constraint $c \in \mathcal{C}_a$}
    \If{$\neg\mathrm{Satisfies}(D_{\mathrm{deploy}}^a, c)$}
      \State $\mathrm{RecordViolation}(c)$
      \If{$c \in \mathcal{C}_{\mathrm{err}}$}
        \State $V_{\mathrm{val}}^a \leftarrow \mathrm{Invalid}$; \ $F \leftarrow F \cup \mathrm{Exclusions}(c)$
      \EndIf
    \EndIf
  \EndFor
\EndFor
\For{each joint constraint $c$ on shared nodes and links}
  \If{$\neg\mathrm{Satisfies}(\{D_{\mathrm{deploy}}^a\}, c)$}
    \State mark the applications involved as Invalid; \ $\mathrm{RecordViolation}(c)$
  \EndIf
\EndFor
\State \Return $\{V_{\mathrm{val}}^a\}$, $F$ \Comment{$F$ is returned to the Deployment Planner}
\end{algorithmic}
\end{algorithm}
